\chapter{Complex Numbers and Roots of Unity}\label{ch:complex}

\begin{center}\itshape
Why is $\cos\frac{2\pi}5=\frac{\sqrt5-1}4$, and why does $\sin\frac\pi7\,\sin\frac{2\pi}7\cdots\sin\frac{6\pi}7=\frac7{64}$?\\
Both answers live on the unit circle.
\end{center}

Complex numbers were not invented because someone wanted $\sqrt{-1}$. They were forced on mathematicians, against their will, by equations
with perfectly real answers. In 1545 Girolamo Cardano published the \emph{Ars Magna}, the book that gave the world the solution of the cubic
equation (and a famous quarrel: Cardano had sworn to Tartaglia in 1539 never to publish the method, and published it when he found that Scipione
del Ferro had discovered it first). In Chapter~37 he asks for two numbers with sum $10$ and product $40$, finds $5+\sqrt{-15}$ and $5-\sqrt{-15}$,
and checks, ``dismissing the mental tortures involved'' (\emph{dimissis incruciationibus}), that their product is indeed $25-(-15)=40$. He then calls
the whole thing ``as subtle as it is useless'' and moves on.

The person who could not move on was Rafael Bombelli, an engineer who drained marshes for a living. Cardano's formula for $x^3=px+q$ reads
\[
x=\sqrt[3]{\frac q2+\sqrt{\frac{q^2}4-\frac{p^3}{27}}}+\sqrt[3]{\frac q2-\sqrt{\frac{q^2}4-\frac{p^3}{27}}},
\]
and for $x^3=15x+4$, whose root $x=4$ is plain to see, it produces $\sqrt[3]{2+\sqrt{-121}}+\sqrt[3]{2-\sqrt{-121}}$. In his \emph{Algebra} of 1572
Bombelli wrote that the idea ``appeared to me more sophistic than true'', and then had the nerve to compute with these impossible quantities as if
they were numbers, with rules he laid down for ``plus of minus'' and ``minus of minus''. He found $(2+\sqrt{-1})^3=2+\sqrt{-121}$, hence
$(2-\sqrt{-1})^3=2-\sqrt{-121}$, and the formula gave $(2+\sqrt{-1})+(2-\sqrt{-1})=4$. The impossible numbers had cancelled out, and the answer was right.

It took two more centuries for the suspicion to fade. Euler gave $\sqrt{-1}$ the name $i$ in 1777. Caspar Wessel, a Norwegian surveyor, presented to
the Danish Academy on 10 March 1797 a paper that drew complex numbers as arrows in the plane and showed that multiplying them means turning and
stretching; it was rediscovered in 1895 and translated into French in 1897, and the picture is named after Jean-Robert Argand, who found it again in 1806. Gauss, who had used the plane
in his 1799 proof that every polynomial has a root, coined the name ``complex number'' in 1831. By then he had long known the most spectacular
consequence: on 30 March 1796, aged eighteen, he had written in his diary that the regular $17$-gon can be drawn with ruler and compass,
something nobody had suspected in two thousand years. The secret, in Bombelli's computation and in Gauss's, is that complex numbers turn geometry
into algebra: a rotation becomes a multiplication, and a regular polygon becomes the set of roots of $z^n=1$.

\portraits{cardano}{Gerolamo Cardano (1501--1576). Woodcut, eighteenth century. Public domain, via Wikimedia Commons.}{euler}{Leonhard Euler (1707--1783). Portrait by Jakob Emanuel Handmann, 1753, Kunstmuseum Basel. Public domain, via Wikimedia Commons.}

\section{Complex numbers from scratch}

\begin{definitionx}
A complex number is a pair of real numbers, written $z=x+iy$, with $x=\operatorname{Re}z$ the \emph{real part} and $y=\operatorname{Im}z$ the
\emph{imaginary part}. Addition is coordinatewise and multiplication follows the single rule $i^2=-1$:
\[
(x+iy)(x'+iy')=(xx'-yy')+i(xy'+yx').
\]
The \emph{conjugate} of $z$ is $\bar z=x-iy$ and its \emph{modulus} is $|z|=\sqrt{x^2+y^2}=\sqrt{z\bar z}$.
\end{definitionx}

One checks directly that addition and multiplication are commutative and associative, that multiplication distributes over addition, and that
every $z\ne0$ has an inverse,
\[
\frac1z=\frac{\bar z}{|z|^2}=\frac{x-iy}{x^2+y^2},
\]
so that $\C$ is a field containing $\R$ (the numbers with $y=0$). We write $\C^*=\C\setminus\{0\}$. Conjugation respects everything: $\overline{z+w}=\bar z+\bar w$,
$\overline{zw}=\bar z\,\bar w$, and $z$ is real if and only if $\bar z=z$. Hence $|zw|^2=zw\,\bar z\bar w=|z|^2|w|^2$: the modulus is multiplicative.

The pair $(x,y)$ is a point of the plane, and $|z|$ is its distance to the origin; $|z-w|$ is the distance between the points $z$ and $w$.
Addition is the addition of vectors. The geometric meaning of multiplication is the subject of the next section; for now, the inequality that
governs all estimates in $\C$:

\begin{propositionx}[triangle inequality]
For $z,w\in\C$, $|z+w|\le|z|+|w|$, with equality if and only if one of them is a non-negative real multiple of the other. Consequently
$\bigl||z|-|w|\bigr|\le|z-w|$.
\end{propositionx}
\begin{proof}
$|z+w|^2=(z+w)(\bar z+\bar w)=|z|^2+|w|^2+z\bar w+\bar zw=|z|^2+|w|^2+2\operatorname{Re}(z\bar w)$, and $\operatorname{Re}(z\bar w)\le|z\bar w|=|z||w|$
because the real part of a number is at most its modulus. So $|z+w|^2\le(|z|+|w|)^2$. Equality needs $\operatorname{Re}(z\bar w)=|z\bar w|$, i.e.\ $z\bar w$
real and $\ge0$. If $w=0$, then $w=0\cdot z$. If $w\ne0$, the condition says $z=\lambda w$ with $\lambda=z\bar w/|w|^2\ge0$. Conversely, if $z=\lambda w$ or $w=\lambda z$ with $\lambda\ge0$, then $|z+w|=(\lambda+1)$ times the modulus of the other factor, which is $|z|+|w|$. The consequence follows from $|z|=|(z-w)+w|\le|z-w|+|w|$ and the same with $z,w$ swapped.
\end{proof}

\begin{examplex}[square roots by algebra]
Does $i$ have a square root? We want $(x+iy)^2=i$, that is $x^2-y^2=0$ and $2xy=1$. The first gives $y=\pm x$; the second then forces $y=x$ and
$x^2=\frac12$. So $\sqrt i$ is $\pm\frac{1+i}{\sqrt2}$. The same two equations, $x^2-y^2=a$ and $2xy=b$, together with $x^2+y^2=|a+ib|$ (take moduli),
give the square roots of any $a+ib$: $x^2=\frac{|a+ib|+a}2$, $y^2=\frac{|a+ib|-a}2$. If $b>0$, then $x$ and $y$ have the same sign; if $b<0$, they have opposite signs; if $b=0$, one of them is $0$ and the other is $\pm\sqrt{|a|}$.
For instance $\sqrt{3+4i}=\pm(2+i)$. Hence every quadratic equation $az^2+bz+c=0$ with complex coefficients and $a\ne0$ has the solutions
$z=\frac{-b\pm\delta}{2a}$ where $\delta^2=b^2-4ac$: complete the square exactly as over $\R$, and the only new fact needed is that every complex
number has a square root. Chapter~\ref{ch:polynomials} proves much more: every nonconstant complex polynomial has a root.
\end{examplex}

\begin{examplex}[Bombelli's computation, checked]
$(2+i)^2=4+4i+i^2=3+4i$, and $(2+i)^3=(3+4i)(2+i)=6+3i+8i+4i^2=2+11i$ (Figure~\ref{fig:bombelli}). Since $\sqrt{-121}=11i$, this is Bombelli's
$\sqrt[3]{2+\sqrt{-121}}=2+\sqrt{-1}$. Exercise~\ref{ex:cardano} derives Cardano's formula and finishes the story: the other two roots of
$x^3=15x+4$ are $-2\pm\sqrt3$, and they too come out of the formula, once one knows that $2+11i$ has three cube roots.
\end{examplex}

\begin{figure}[ht]
\centering
\fig{ch02-bombelli}
\caption{The powers $1$, $2+i$, $3+4i$, $2+11i$ of $2+i$. Each multiplication by $2+i$ multiplies the modulus by $\sqrt5$ and turns the arrow by the same angle, about $26.6^\circ$: the points lie on a logarithmic spiral. The exponential form of the next section explains why.}
\label{fig:bombelli}
\end{figure}

\begin{think}
Multiply $1+i$ by itself: you get $2i$. Draw $1+i$ and $2i$ as arrows. What happened to the length, and what happened to the angle with the real axis? Now multiply $1+i$ by $i$: the arrow $-1+i$. Guess the general rule before reading on.
\end{think}

% Complex series, after the arithmetic of the modulus. The real tests are in the sequences chapter.
\subsection*{Convergent series with complex terms}
A complex sequence $z_n=x_n+iy_n$ converges to $x+iy$ exactly when
$x_n\to x$ and $y_n\to y$: use $|x_n-x|,|y_n-y|\le|z_n-z|$ and
$|z_n-z|\le|x_n-x|+|y_n-y|$. Thus Cauchy's criterion, absolute
convergence, the Cauchy product and Dirichlet's test from
Chapter~\ref{ch:sequences} extend to complex sequences. For the Cauchy
criterion apply the real result to each coordinate; the other proofs are
unchanged, using the complex triangle inequality.

\section{The exponential form}

Use the real functions $\exp$, $\cos$ and $\sin$, and the constant $\pi$, with the elementary real properties listed in the foundational contract. For $z=x+iy$, define
\[
\exp z=e^x(\cos y+i\sin y).
\]
This agrees with the real exponential when $y=0$. The real exponential and trigonometric addition laws, followed by ordinary complex multiplication, give $\exp(z+w)=\exp z\exp w$. Evenness of $\cos$ and oddness of $\sin$ give $\overline{\exp z}=\exp\bar z$. We shall write $e^z$ for $\exp z$. The absolutely convergent series $\sum_{n\ge0}z^n/n!$ agrees with this definition; its identification is proved with Taylor estimates in Exercise~\ref{ex:taylorln}(c), a return topic after Chapter~\ref{ch:calculus}. It is not an assumption needed for the algebra here.

\begin{proposition}[exponential form]\label{prop:expform}
\begin{enumerate}[label=(\alph*)]
\item $|e^{i\theta}|=1$, $e^{i\alpha}e^{i\beta}=e^{i(\alpha+\beta)}$, and hence (de Moivre) $(\cos\theta+i\sin\theta)^n=\cos n\theta+i\sin n\theta$ for $n\in\Z$.
\item Every $z\in\C^*$ can be written $z=r\,e^{i\theta}$ with $r=|z|$ and $\theta\in\R$, unique modulo $2\pi$.
\item (Euler) $\cos\theta=\dfrac{e^{i\theta}+e^{-i\theta}}2$ and $\sin\theta=\dfrac{e^{i\theta}-e^{-i\theta}}{2i}$.
\end{enumerate}
\end{proposition}
\begin{proof}
(a) By conjugation and the addition law just proved, $|e^{i\theta}|^2=e^{i\theta}e^{-i\theta}=e^0=1$. The product rule is $\exp(a+b)=\exp a\exp b$, and de Moivre follows by induction (and $e^{-i\theta}=(e^{i\theta})^{-1}$ for negative $n$).

(c) $e^{-i\theta}=\overline{e^{i\theta}}=\cos\theta-i\sin\theta$; add and subtract.

(b) Put $r=|z|$ and $u=z/r=x+iy$, so $x^2+y^2=1$. As $\cos$ maps $[0,\pi]$ onto $[-1,1]$, there is $\theta_0\in[0,\pi]$ with $\cos\theta_0=x$; then $\sin^2\theta_0=1-x^2=y^2$. Proposition~\ref{prop:sincmp} gives $\sin\ge0$ on $[0,\pi]$, so $\sin\theta_0=|y|$, and $\theta=\pm\theta_0$ gives $u=e^{i\theta}$. For uniqueness, $e^{i\theta}=e^{i\theta'}$ means $e^{i(\theta-\theta')}=1$, i.e.\ $\cos(\theta-\theta')=1$ and $\sin(\theta-\theta')=0$. Write $\theta-\theta'=2\pi k+t$ with $k\in\Z$ and $t\in[0,2\pi[$. Since $\cos\pi=-1$ and $|e^{i\pi}|=1$, $e^{i\pi}=-1$ and $e^{2i\pi}=1$, so $e^{it}=e^{i(\theta-\theta')}=1$ and $\cos t=1$. On $[0,\pi]$, $\cos$ is injective, so $\cos t=1$ forces $t=0$ there; on $]\pi,2\pi[$, $\cos t=\cos(2\pi-t)$ with $2\pi-t\in]0,\pi[$, where $\cos<1$. So $t=0$ and $\theta-\theta'\in2\pi\Z$.
\end{proof}

The number $\theta$ is an \emph{argument} of $z$. Statement (b) is the geometry of multiplication, promised above: $re^{i\theta}\cdot se^{i\varphi}=rs\,e^{i(\theta+\varphi)}$,
so multiplying by $w=se^{i\varphi}$ stretches every arrow by $s$ and turns it by the angle $\varphi$ (Figure~\ref{fig:multiply}). Multiplying by $i=e^{i\pi/2}$ is a
quarter turn, which is why $i^2=-1$: two quarter turns make a half turn. Bombelli's spiral in Figure~\ref{fig:bombelli} is the picture of $(\sqrt5e^{i\theta})^n$
with $\theta=\arctan\frac12$.

\begin{figure}[ht]
\centering
\fig{ch02-multiply}
\caption{Multiplication by $w=1.2\,e^{i\cdot50^\circ}$: the arrow $z=2+i$ is stretched by $1.2$ and turned by $50^\circ$. In general $|zw|=|z|\,|w|$ and the arguments add.}
\label{fig:multiply}
\end{figure}

Two reflexes follow. To \emph{linearise} a power of $\cos$ or $\sin$, expand with Euler and the binomial theorem.
To \emph{factor} a sum like $e^{ia}+e^{ib}$, pull out the half-angle:
\[
\begin{aligned}e^{ia}+e^{ib}&=e^{i\frac{a+b}2}\left(e^{i\frac{a-b}2}+e^{-i\frac{a-b}2}\right)\eqbr=2\cos\frac{a-b}2\;e^{i\frac{a+b}2}.\end{aligned}
\]
In particular $1+e^{i\theta}=2\cos\frac\theta2\,e^{i\theta/2}$ and $1-e^{i\theta}=-2i\sin\frac\theta2\,e^{i\theta/2}$,
so $|1-e^{i\theta}|=2\left|\sin\frac\theta2\right|$. The half-angle trick is worth a second look, because it will be used a dozen times: $e^{ia}+e^{ib}$ is
the sum of two unit arrows, and the sum of two arrows of equal length points along the bisector, at angle $\frac{a+b}2$; its length is $2\left|\cos\frac{a-b}2\right|$ (if the cosine is negative, the resultant points opposite to the chosen half-angle direction) by
the geometry of the rhombus. The algebra above is that picture.

\begin{example}
\[\begin{aligned}\cos^3\theta&=\frac18\left(e^{i\theta}+e^{-i\theta}\right)^3\eqbr=\frac18\left(e^{3i\theta}+3e^{i\theta}+3e^{-i\theta}+e^{-3i\theta}\right)
\eqbr=\frac{\cos3\theta+3\cos\theta}4.\end{aligned}\]
\end{example}

\begin{examplex}[the addition formulas, encoded]
Take real and imaginary parts of $e^{i(\alpha+\beta)}=e^{i\alpha}e^{i\beta}=(\cos\alpha+i\sin\alpha)(\cos\beta+i\sin\beta)$:
\[
\begin{gathered}\cos(\alpha+\beta)=\cos\alpha\cos\beta-\sin\alpha\sin\beta,\qq\sin(\alpha+\beta)=\sin\alpha\cos\beta+\cos\alpha\sin\beta.\end{gathered}
\]
The exponential notation packages the real addition formulas used in its construction. Recovering them here is a useful consistency check, rather than an independent proof of those admitted real facts. Many further identities now follow by ordinary algebra. The sum-to-product formulas are the half-angle trick; the double-angle formulas are $\alpha=\beta$; $\cos2\theta=2\cos^2\theta-1$
is the real part of $(e^{i\theta})^2$ written as $(\cos\theta+i\sin\theta)^2$, with $\sin^2=1-\cos^2$.
\end{examplex}

\section{Roots of unity}

\begin{think}
Draw the solutions of $z^6=1$ in the plane. Which of them have all six as powers, and which only some? What is the sum of all six? Answer from the picture before computing anything.
\end{think}

\begin{proposition}[roots of unity]\label{prop:roots}
Let $n\ge1$ and $\omega=e^{2i\pi/n}$.
\begin{enumerate}[label=(\alph*)]
\item The equation $z^n=1$ has exactly $n$ solutions, the $\omega^k$ for $0\le k\le n-1$. They form a group $\mathbb U_n$ under multiplication, and
$X^n-1=\prod_{k=0}^{n-1}\left(X-\omega^k\right)$.
\item The order of $\omega^k$ is $\frac n{\gcd(k,n)}$. In particular $\omega^k$ is \emph{primitive} (it generates $\mathbb U_n$) exactly when $\gcd(k,n)=1$.
\end{enumerate}
\end{proposition}
\begin{proof}
(a) If $z^n=1$ then $|z|^n=1$, so $|z|=1$ and $z=e^{i\theta}$; then $e^{in\theta}=1$, so $n\theta\in2\pi\Z$ by uniqueness in Proposition~\ref{prop:expform}, i.e.\ $z=\omega^k$ for some $k\in\Z$, and $k$ only matters modulo $n$. Conversely each $\omega^k$ is a solution, and $\omega^0,\dots,\omega^{n-1}$ are distinct (their arguments $\frac{2k\pi}n$ are distinct modulo $2\pi$). Products and inverses of solutions are solutions. Finally, the elementary identity $X^j-a^j=(X-a)\sum_{l=0}^{j-1}X^{j-1-l}a^l$ shows that $P(X)-P(a)$ is divisible by $X-a$. Thus a root $a$ gives a factor $X-a$. Removing one factor leaves every other distinct root, since $0=(b-a)Q(b)$ and $b\ne a$. Repeating gives divisibility by the product of all $n$ factors; monicity and degree give equality. Proposition~\ref{prop:rootsfactor} develops this argument systematically.

(b) Let $d=\gcd(k,n)$. Then $(\omega^k)^m=1$ iff $n\mid km$ iff $\frac nd\mid\frac kd m$ iff $\frac nd\mid m$, the last step by Gauss' lemma since $\gcd\!\left(\frac nd,\frac kd\right)=1$. So the order is $\frac nd$, and $\omega^k$ generates the $n$-element group $\mathbb U_n$ iff its order is $n$, iff $d=1$.
\end{proof}

The $n$-th roots of unity are the vertices of the regular $n$-gon inscribed in the unit circle, with one vertex at $1$. The ``order'' of a root is the
number of its distinct powers; a primitive root visits every vertex before returning to $1$, a non-primitive one only visits a smaller regular polygon
(Figure~\ref{fig:roots}, left). For $n=6$: $\omega$ and $\omega^5$ are primitive, $\omega^2$ and $\omega^4$ have order $3$ (they are the cube roots of unity),
$\omega^3=-1$ has order $2$.

\begin{figure}[ht]
\centering
\fig{ch02-roots}
\caption{Left: the twelve $12$th roots of unity. The primitive ones, $\omega^k$ with $\gcd(k,12)=1$, are in red; the others generate the smaller groups $\mathbb U_1,\mathbb U_2,\mathbb U_3,\mathbb U_4,\mathbb U_6$. Right: the six chords from $1$ in the regular heptagon have lengths $|1-\omega^k|=2\sin\frac{k\pi}7$, and their product is exactly $7$ (Exercise~\ref{ex:sines}).}
\label{fig:roots}
\end{figure}

\begin{think}
Inscribe a regular heptagon in the unit circle and multiply the lengths of the six chords from one vertex to the others (Figure~\ref{fig:roots}, right). A calculator gives the answer: it is an integer. Which one, and what would you guess for an $n$-gon? The polynomial $X^n-1$ holds the explanation.
\end{think}

\begin{proposition}[orthogonality]\label{prop:orth}
Let $\omega=e^{2i\pi/n}$ and $j\in\Z$. Then
\[
\sum_{k=0}^{n-1}\omega^{jk}=\begin{cases}n&\text{if }n\mid j,\\0&\text{otherwise.}\end{cases}
\]
\end{proposition}
\begin{proof}
If $n\mid j$ every term is $1$. Otherwise $q=\omega^j\ne1$ and $q^n=1$, so the geometric sum is $\frac{q^n-1}{q-1}=0$.
\end{proof}

The case $j=1$ says that the vertices of a regular polygon have centre of mass at the origin: the arrows $1,\omega,\dots,\omega^{n-1}$ cancel exactly
(Figure~\ref{fig:filter}). The name ``orthogonality'' comes from the following reading: the vectors $v_j=(1,\omega^j,\omega^{2j},\dots,\omega^{(n-1)j})$, for
$j=0,\dots,n-1$, are pairwise orthogonal for the Hermitian product $\langle v,w\rangle=\sum_kv_k\overline{w_k}$, since $\langle v_j,v_l\rangle=\sum_k\omega^{(j-l)k}$ is $0$ for $j\ne l$.
These $v_j$ are the \emph{Fourier basis}, and Exercise~\ref{ex:circulant} shows what they are good for.

\section{The roots-of-unity filter}

\begin{think}
How do you get $\sum_k\binom n{2k}$ (the even-index coefficients of $(1+X)^n$) from $(1+1)^n$ and $(1-1)^n$? Now you want only the indices divisible by $3$: what should replace $1$ and $-1$?
\end{think}

Proposition~\ref{prop:orth} lets us extract from a polynomial only the terms whose exponent lies in a given residue class.

\begin{theorem}[roots-of-unity filter]\label{thm:filter}
Let $P(X)=\sum_k a_kX^k$, let $m\ge1$, $\omega=e^{2i\pi/m}$ and $r\in\Z$. Then
\[
\sum_{k\equiv r\ (\mathrm{mod}\ m)}a_k=\frac1m\sum_{j=0}^{m-1}\omega^{-jr}\,P\!\left(\omega^j\right).
\]
\end{theorem}
\begin{proof}
The right-hand side is
\[
\sum_ka_k\cdot\frac1m\sum_{j=0}^{m-1}\omega^{j(k-r)},
\]
and by Proposition~\ref{prop:orth} the inner average is $1$ if $m\mid k-r$ and $0$ otherwise.
\end{proof}

\begin{figure}[ht]
\centering
\fig{ch02-filter}
\caption{The filter with $m=3$, seen as arrows. For an exponent $k$ not divisible by $3$, the three numbers $1,\omega^k,\omega^{2k}$ are the vertices of an equilateral triangle and add up to $0$; for $3\mid k$ they all equal $1$. Averaging $P(\omega^j)$ over $j$ therefore keeps exactly the coefficients $a_k$ with $3\mid k$.}
\label{fig:filter}
\end{figure}

\begin{example}
With $P=(1+X)^n$, $m=2$, $r=0$: $\sum_k\binom n{2k}=\frac12\left(2^n+0^n\right)=2^{n-1}$ for $n\ge1$.
The case $m=3$ is Exercise~\ref{ex:mod3}.
\end{example}

\section{Geometry with complex numbers}

Every map of the form $z\mapsto az+b$ with $a\ne0$ is a \emph{direct similarity}: writing $a=re^{i\theta}$, it is a rotation by $\theta$ followed by a
scaling by $r$ and a translation by $b$. Such maps preserve angles and ratios of lengths, so they send equilateral triangles to equilateral triangles,
squares to squares, and so on. Conversely the only direct similarities of the plane are these (the indirect similarities are $z\mapsto a\bar z+b$ with $a\ne0$; reflections are the special cases with $|a|=1$ and $a\bar b+b=0$), a fact we do not need.
The practical rule: a geometric statement that is invariant under similarities may be proved after moving two convenient points to $0$ and $1$.
Here is the criterion that handles most problems about equilateral triangles. Write $j=e^{2i\pi/3}=\frac{-1+i\sqrt3}2$; it satisfies $j^3=1$, $1+j+j^2=0$, $\bar j=j^2$,
and $-j^2=e^{i\pi/3}$ is the rotation by $60^\circ$.

\begin{proposition}[equilateral triangles]\label{prop:equilateral}
For $a,b,c\in\C$, the following are equivalent:
\begin{enumerate}[label=(\roman*)]
\item the triangle $abc$ is equilateral (possibly degenerate: $a=b=c$);
\item $a+jb+j^2c=0$ or $a+j^2b+jc=0$;
\item $a^2+b^2+c^2=ab+bc+ca$.
\end{enumerate}
The first alternative in (ii) holds when $a,b,c$ are in counterclockwise order, the second when they are clockwise.
\end{proposition}
\begin{proof}
The triangle is equilateral and counterclockwise exactly when $c$ is obtained from $b$ by the rotation of angle $60^\circ$ about $a$:
$c-a=-j^2(b-a)$, that is $c-a+j^2b-j^2a=0$, i.e.\ $-a(1+j^2)+j^2b+c=0$, and since $-(1+j^2)=j$ this reads $ja+j^2b+c=0$; multiplying by $j^2$ gives $a+jb+j^2c=0$.
The clockwise case is the rotation by $-60^\circ$, $c-a=-j(b-a)$, which the same computation turns into $a+j^2b+jc=0$. For (iii), expand the product
\[
(a+jb+j^2c)(a+j^2b+jc)=a^2+b^2+c^2-ab-bc-ca,
\]
using $j^3=1$ and $j+j^2=-1$ (the cross terms are $ab(j^2+j)$, $bc(j^2+j)$ and $ca(j+j^2)$). It vanishes exactly when one of the factors does.
\end{proof}

\begin{figure}[ht]
\centering
\fig{ch02-napoleon}
\caption{Napoleon's theorem (Exercise~\ref{ex:napoleon}): on the sides of any triangle $abc$, build equilateral triangles pointing outward; their centres (green dots) form an equilateral triangle (dashed), whatever the shape of $abc$.}
\label{fig:napoleon}
\end{figure}

\section{Cyclotomic polynomials}

The polynomial $X^n-1$ factors over $\C$ into the $n$ linear factors $X-\omega^k$. Grouping the roots by their order gives a factorisation over $\Z$
that is the key to everything the series will say about roots of unity, from the $17$-gon (Volume~2) to primes in arithmetic progressions (Volume~5).

\begin{definition}\label{def:cyclotomic}
The $n$-th \emph{cyclotomic polynomial} is $\Phi_n(X)=\prod_\zeta(X-\zeta)$, the product over the primitive $n$-th roots of unity, i.e.\ over
$\zeta=\omega^k$ with $1\le k\le n$ and $\gcd(k,n)=1$. Its degree is the number of such $k$, which Chapter~\ref{ch:counting} (Exercise~\ref{ex:phi}) computes: it is Euler's $\varphi(n)$.
\end{definition}

\begin{proposition}\label{prop:cyclotomic}
\begin{enumerate}[label=(\alph*)]
\item $X^n-1=\prod_{d\mid n}\Phi_d(X)$, the product over the positive divisors $d$ of $n$.
\item $\Phi_n$ is monic with integer coefficients.
\item For $p$ prime, $\Phi_p=1+X+\dots+X^{p-1}$, and $\Phi_{p^k}(X)=\Phi_p\bigl(X^{p^{k-1}}\bigr)$.
\end{enumerate}
\end{proposition}
\begin{proof}
(a) Every $n$-th root of unity has an order $d$ which divides $n$ (Proposition~\ref{prop:roots}(b)), and the roots of order exactly $d$ are the primitive $d$-th roots of
unity, which are $n$-th roots of unity because $d\mid n$. So the roots of $X^n-1$, each simple, are partitioned according to their order, and
$X^n-1=\prod_k(X-\omega^k)=\prod_{d\mid n}\prod_{\operatorname{ord}\zeta=d}(X-\zeta)=\prod_{d\mid n}\Phi_d$.

(b) Strong induction on $n$. $\Phi_1=X-1$. For $n>1$, (a) gives $X^n-1=\Phi_n\cdot Q$ with $Q=\prod_{d\mid n,\,d<n}\Phi_d$ monic with integer
coefficients by induction. Cancel the leading term of $X^n-1$ by subtracting a suitable integer multiple of a power of $X$ times $Q$, and repeat until the remainder has degree below $\deg Q$. All coefficients stay integral because $Q$ is monic. If two quotient--remainder pairs existed, their difference would give a nonzero multiple of $Q$ of degree below $\deg Q$, impossible. Since (a) already gives zero remainder over $\C$, the integer quotient is $\Phi_n$. This is monic polynomial division; Proposition~\ref{prop:polydiv} gives the general algorithm.

(c) All $p$-th roots of unity but $1$ are primitive, so $\Phi_p=\frac{X^p-1}{X-1}$. For $p^k$, the non-primitive $p^k$-th roots are those of order dividing $p^{k-1}$,
i.e.\ the roots of $X^{p^{k-1}}-1$, so $\Phi_{p^k}=\frac{X^{p^k}-1}{X^{p^{k-1}}-1}=\Phi_p(X^{p^{k-1}})$.
\end{proof}

\begin{examplex}
$\Phi_1=X-1$, $\Phi_2=X+1$, $\Phi_3=X^2+X+1$, $\Phi_4=X^2+1$, $\Phi_6=X^2-X+1$ (from $X^6-1=\Phi_1\Phi_2\Phi_3\Phi_6$), $\Phi_8=X^4+1$,
$\Phi_{12}=X^4-X^2+1$. For $n\le104$ all coefficients are $0$ or $\pm1$, and it is tempting to conjecture that this always holds; $\Phi_{105}$ has two
coefficients equal to $-2$ (Migotti, 1883), and Schur showed that the coefficients are unbounded. Exercise~\ref{ex:cyclo} computes a few more, and Chapter~\ref{ch:polynomials} proves
that $\Phi_p$ is irreducible over $\Q$; the irreducibility of every $\Phi_n$ is a theorem of Gauss, proved in Volume~2.
\end{examplex}

\portrait{demoivre}{Abraham de Moivre (1667--1754). Portrait, 1736, artist unknown. Public domain, via Wikimedia Commons.}

\section{Workshop: complex numbers are a steering wheel}
A complex number has two jobs: a position when added, and a turn-and-stretch instruction when multiplied. Confusing these jobs is like trying to drive by adding the steering wheel to the road.

\begin{examplex}[a robot that nearly returns home]
A robot takes successive steps $1,q,q^2,q^3$, starting at $0$. If $q=i$, its successive positions are $0,1,1+i,i,0$: a square and a return home. If $q=(1+i)/\sqrt2$, eight unit steps return home, by the geometric sum. But if $q=1+i$, each step is $\sqrt2$ times the preceding length. After four steps its position is
\[
1+(1+i)+2i+(-2+2i)=5i,
\]
not $0$. An angle calculation alone forgot the lengths. In general after $n$ steps, the position is $(1-q^n)/(1-q)$ for $q\ne1$; for $q=1$ it is $n$. Separate that exceptional value before dividing.
\end{examplex}

\begin{figure}[ht]\centering\fig{review02-robot}
\caption{The robot with step multiplier $1+i$. Successive positions are $0,1,2+i,2+3i,5i$; arrows are its four steps, not its position vectors. The turning angle is always $45^\circ$, but the lengths grow.}\end{figure}

\subsection*{Three doors in the argument department}
\textbf{Mission.} Solve $z^3=8i$ completely, then explain why the principal argument of a product can require a correction of $2\pi$.

\textbf{Hint.} Write $8i=8e^{i\pi/2}$; divide \emph{all} possible arguments by three, not just the first one.

\textbf{Solution.} Every solution has modulus $2$ and argument $\pi/6+2k\pi/3$ for $k=0,1,2$. Thus the roots are $\sqrt3+i$, $-\sqrt3+i$ and $-2i$. Cubing checks them directly, and their distinct arguments prove they are distinct. There are no others by the exponential form. If principal arguments are chosen in $]-\pi,\pi]$, two numbers of argument $3\pi/4$ have product of argument $3\pi/2$, represented principally by $-\pi/2$. Arguments add modulo $2\pi$, not always as chosen representatives.

\begin{reminder}[a square-root symbol needs a convention]
The equation $w^2=z\ne0$ has two solutions, but a chosen square-root \emph{function} returns just one. In particular the real identity $\sqrt{ab}=\sqrt a\sqrt b$ cannot be used indiscriminately for complex principal roots: for $a=b=-1$, the two sides are $1$ and $-1$. An equation with all its roots is safer than an unannounced choice of branch.
\end{reminder}


\section{Exercises}

\begin{exercise}[$\star$]\label{ex:lin}
Linearise $\sin^4\theta$, and, after the integration tools of Chapter~\ref{ch:calculus}, deduce $\displaystyle\int_0^{\pi}\sin^4\theta\dd\theta$.
\end{exercise}

\begin{exercise}[$\star$]\label{ex:binomcos}
Compute $\displaystyle\sum_{k=0}^n\binom nk\cos(k\theta)$ in closed form.
\end{exercise}

\begin{exercise}[$\star\star$]\label{ex:dirichlet}
\emph{The Dirichlet kernel.} For $\theta\notin2\pi\Z$, show that
\[
D_n(\theta)=\sum_{k=-n}^ne^{ik\theta}=\frac{\sin\left((n+\frac12)\theta\right)}{\sin\frac\theta2},
\]
and deduce $\displaystyle\sum_{k=1}^n\cos k\theta$.
\end{exercise}

\begin{exercise}[$\star\star$]\label{ex:mod3}
\emph{Bridge: Binomials and roots of unity.}
How many subsets of $\{1,\dots,n\}$ have a number of elements divisible by $3$?
\end{exercise}

\begin{exercise}[$\star\star$]\label{ex:sines}
Show that $\displaystyle\prod_{k=1}^{n-1}\sin\frac{k\pi}n=\frac n{2^{n-1}}$ for $n\ge2$.
\emph{Hint: evaluate $\prod_{k=1}^{n-1}(X-\omega^k)$ at $X=1$.}
\end{exercise}

\begin{exercise}[$\star\star$]\label{ex:pentagon}
Let $\omega=e^{2i\pi/5}$, $c=\omega+\omega^4$ and $d=\omega^2+\omega^3$. Compute $c+d$ and $cd$, and deduce that
$\cos\frac{2\pi}5=\frac{\sqrt5-1}4$. (This is why the regular pentagon can be drawn with ruler and compass.)
\end{exercise}

\begin{exercise}[$\star\star\star$]\label{ex:circulant}
\emph{Bridge: Roots of unity and linear algebra.}
\emph{Circulant determinants (return after Chapter~\ref{ch:linear}).} Let $a_0,\dots,a_{n-1}\in\C$, $P(X)=\sum_ka_kX^k$, and let $C$ be the $n\times n$ matrix
with entries $C_{r,s}=a_{(s-r)\bmod n}$ (each row is the previous one shifted right by one).
\begin{enumerate}[label=(\alph*)]
\item With $\omega=e^{2i\pi/n}$ and $v_j=\left(1,\omega^j,\omega^{2j},\dots,\omega^{(n-1)j}\right)^{\mathsf T}$, show that $Cv_j=P(\omega^j)\,v_j$.
\item Deduce that $\det C=\prod_{j=0}^{n-1}P(\omega^j)$.
\item For $n\ge2$, compute the determinant of the circulant matrix with first row $(1,1,0,\dots,0)$.
\end{enumerate}
\end{exercise}

\begin{exercise}[$\star\star\star$]\label{ex:gauss}
\emph{A quadratic Gauss sum.} Let $p$ be an odd prime, $\omega=e^{2i\pi/p}$ and $G=\sum_{a=0}^{p-1}\omega^{a^2}$. Show that $|G|^2=p$.
\end{exercise}

\subsection*{Further exercises}

\begin{exercise}[$\star$]\label{ex:sqrtz}
Find the square roots of $-7+24i$ and of $2i$, and solve $z^2-(1+i)z+i=0$.
\end{exercise}

\begin{exercise}[$\star$]\label{ex:sumdist}
Let $n\ge2$, let $z_0,\dots,z_{n-1}$ be the vertices of a regular $n$-gon inscribed in the unit circle, and $z$ any point of that circle. Show that
$\sum_{k=0}^{n-1}|z-z_k|^2=2n$.
\end{exercise}

\begin{exercise}[$\star\star$]\label{ex:cardano}
\emph{Cardano's formula, and Bombelli's rescue.}
\begin{enumerate}[label=(\alph*)]
\item Let $p,q\in\R$. Show that if $u,v$ satisfy $u^3+v^3=q$ and $uv=\frac p3$, then $x=u+v$ solves $x^3=px+q$. Deduce that $u^3$ and $v^3$ are the roots of $T^2-qT+\frac{p^3}{27}$, hence Cardano's formula.
\item For $x^3=15x+4$, show that $u^3=2+11i$, $v^3=2-11i$, and find all three cube roots of $2+11i$ in the form $(2+i)\omega^k$ with $\omega=e^{2i\pi/3}$. Pair each with the $v$ for which $uv=5$, and recover the three roots $4$, $-2+\sqrt3$, $-2-\sqrt3$.
\end{enumerate}
\end{exercise}

\begin{exercise}[$\star\star$]\label{ex:equi}
\begin{enumerate}[label=(\alph*)]
\item Let $a,b,c$ be three distinct complex numbers with $|a|=|b|=|c|=1$ and $a+b+c=0$. Show that $abc$ is an equilateral triangle.
\item Let $abc$ be equilateral. Show that the triangle with vertices $a+b$, $b+c$, $c+a$ is equilateral.
\end{enumerate}
\end{exercise}

\begin{exercise}[$\star\star$]\label{ex:cyclo}
\begin{enumerate}[label=(\alph*)]
\item Compute $\Phi_{10}$, $\Phi_{15}$ and $\Phi_{18}$ from Proposition~\ref{prop:cyclotomic}(a).
\item Show that $\Phi_{2n}(X)=\Phi_n(-X)$ for odd $n\ge3$.
\item Show that $\Phi_n(1)=p$ if $n$ is a power of a prime $p$, and $\Phi_n(1)=1$ if $n$ has at least two distinct prime factors. (Evaluate $\frac{X^n-1}{X-1}=\prod_{d\mid n,\,d>1}\Phi_d$ at $X=1$ and use strong induction.)
\end{enumerate}
\end{exercise}

\begin{exercise}[$\star\star\star$]\label{ex:napoleon}
\emph{Napoleon's theorem.} On the sides of a triangle $abc$, taken counterclockwise, build equilateral triangles pointing outward, and let $A,B,C$ be their
centres ($A$ on the side $bc$, $B$ on $ca$, $C$ on $ab$). Show that $ABC$ is equilateral. \emph{Hint: the outward vertex on the side $bc$ is $b+(c-b)e^{-i\pi/3}$; express $A,B,C$ and apply Proposition~\ref{prop:equilateral}.}
\end{exercise}

\begin{exercise}[$\star\star\star$]\label{ex:cotsum}
\emph{A finite Basel sum (return after formal differentiation in Chapter~\ref{ch:polynomials}).} Let $n\ge2$, $\omega=e^{2i\pi/n}$ and $P(X)=1+X+\dots+X^{n-1}=\prod_{k=1}^{n-1}(X-\omega^k)$.
\begin{enumerate}[label=(\alph*)]
\item Show that $\dfrac{P'}{P}=\displaystyle\sum_{k=1}^{n-1}\frac1{X-\omega^k}$ and $\dfrac{P'^2-P''P}{P^2}=\displaystyle\sum_{k=1}^{n-1}\frac1{(X-\omega^k)^2}$, and evaluate both at $X=1$.
\item Show that $\dfrac1{1-\omega^k}=\dfrac12+\dfrac i2\cot\dfrac{k\pi}n$, and deduce $\displaystyle\sum_{k=1}^{n-1}\cot^2\frac{k\pi}n=\frac{(n-1)(n-2)}3$ and $\displaystyle\sum_{k=1}^{n-1}\frac1{\sin^2\frac{k\pi}n}=\frac{n^2-1}3$.
\item Deduce that $\displaystyle\sum_{k=1}^{n-1}\frac1{|1-\omega^k|^2}=\frac{n^2-1}{12}$, and, using $\sin x\sim x$, guess what this says about $\sum\frac1{k^2}$ as $n\to\infty$. (Chapter~\ref{ch:polynomials} makes the guess a proof.)
\end{enumerate}
\end{exercise}

\begin{challenge}
\begin{enumerate}[label=(\arabic*)]
\item \emph{The sign of the Gauss sum.} Compute $G=\sum_a\omega^{a^2}$ numerically for $p=3,5,7,11,13$. The pattern, $G=\sqrt p$ for $p\equiv1\pmod4$ and $G=i\sqrt p$ for $p\equiv3\pmod4$, was conjectured by Gauss in 1801 and took him four years to prove; Volume~4 gives a proof.
\item \emph{Constructible polygons.} Exercise~\ref{ex:pentagon} expresses $\cos\frac{2\pi}5$ with one square root. Try to express $\cos\frac{2\pi}7$ with square roots alone; you will fail, and Volume~2 explains why: the regular $n$-gon is constructible with ruler and compass exactly when $\varphi(n)$ is a power of $2$ (Gauss 1796, Wantzel 1837). Which $n\le30$ qualify?
\item \emph{Sums of roots of unity.} Find all triples of $n$-th roots of unity $\alpha,\beta,\gamma$ with $\alpha+\beta=\gamma$. (Draw.) Then find all solutions of $\alpha+\beta+\gamma+\delta=0$ among roots of unity.
\end{enumerate}
Notes on the challenges are at the end of the chapter.
\end{challenge}

\begin{beauty}
The filter of Theorem~\ref{thm:filter} is the discrete Fourier transform in disguise: averaging $P$ against the
characters $j\mapsto\omega^{-jr}$ isolates one frequency. Volume~4 turns this into the FFT, an algorithm that multiplies
polynomials (and large integers) fast. The Gauss sum of Exercise~\ref{ex:gauss} returns in Volume~1, where its
\emph{square} $G^2=\pm p$ gives a proof of quadratic reciprocity; its sign is a much deeper theorem of Gauss (Volume~4). Exercise~\ref{ex:circulant} is a first example of a matrix
diagonalised by the Fourier basis, the idea behind the characters of Volume~2. And the cyclotomic polynomials, so innocent here, carry the arithmetic of
roots of unity through the whole series: their irreducibility is Galois theory's first real theorem, and their values at integers are what Euclid's proof
of the infinitude of primes needs to reach the primes $\equiv1\pmod n$.
\end{beauty}

\section{Hints}

\begin{hint}{ex:lin}
Expand $\left(\frac{e^{i\theta}-e^{-i\theta}}{2i}\right)^4$ with the binomial formula and group conjugate terms into cosines.
\end{hint}
\begin{hint}{ex:binomcos}
It is the real part of $\sum_k\binom nke^{ik\theta}=(1+e^{i\theta})^n$. Factor out the half-angle: $1+e^{i\theta}=2\cos\frac\theta2\,e^{i\theta/2}$.
\end{hint}
\begin{hint}{ex:dirichlet}
A geometric sum with ratio $e^{i\theta}$. Multiply numerator and denominator by $e^{-i\theta/2}$ to make sines appear.
\end{hint}
\begin{hint}{ex:mod3}
Apply Theorem~\ref{thm:filter} to $(1+X)^n$ with $m=3$, $r=0$. Write $1+j$ and $1+j^2$ in exponential form.
\end{hint}
\begin{hint}{ex:sines}
$\prod_{k=1}^{n-1}(X-\omega^k)=\frac{X^n-1}{X-1}=1+X+\dots+X^{n-1}$. Evaluate at $1$, take moduli, and use $|1-e^{i\theta}|=2\left|\sin\frac\theta2\right|$.
\end{hint}
\begin{hint}{ex:pentagon}
$c+d=\omega+\omega^2+\omega^3+\omega^4=-1$. For $cd$, expand and reduce the exponents modulo $5$. Then $c$ and $d$ are the roots of a quadratic, and $c=2\cos\frac{2\pi}5>0$.
\end{hint}
\begin{hint}{ex:circulant}
(a) The $r$-th coordinate of $Cv_j$ is $\sum_sa_{(s-r)\bmod n}\omega^{sj}$; substitute $t=s-r$ modulo $n$. (b) The $v_j$ form a basis (a Vandermonde matrix). (c) $P=1+X$.
\end{hint}
\begin{hint}{ex:gauss}
$|G|^2=\sum_{a,b}\omega^{a^2-b^2}$. Substitute $a=b+h$: the exponent becomes $h^2+2bh$, linear in $b$. Then sum over $b$ first with Proposition~\ref{prop:orth}.
\end{hint}
\begin{hint}{ex:sqrtz}
The method of the first section: $x^2-y^2=a$, $2xy=b$, $x^2+y^2=|a+ib|$. For the equation, compute the discriminant and take its square root the same way.
\end{hint}
\begin{hint}{ex:sumdist}
$|z-z_k|^2=(z-z_k)(\bar z-\bar z_k)=2-2\operatorname{Re}(z\bar z_k)$, and $\sum_k\bar z_k=0$.
\end{hint}
\begin{hint}{ex:cardano}
(a) Expand $(u+v)^3$. (b) $|2+11i|=5\sqrt5=(\sqrt5)^3$, so every cube root has modulus $\sqrt5$; $v$ must be $\frac5u$, which for $u=(2+i)\omega^k$ is $(2-i)\omega^{-k}$.
\end{hint}
\begin{hint}{ex:equi}
(a) $a^2+b^2+c^2=(a+b+c)^2-2(ab+bc+ca)$, and $ab+bc+ca=abc(\bar a+\bar b+\bar c)$ on the unit circle. (b) Apply the criterion (ii) to the new vertices.
\end{hint}
\begin{hint}{ex:cyclo}
(a) $X^{10}-1=\Phi_1\Phi_2\Phi_5\Phi_{10}$, and so on: divide. (b) The primitive $2n$-th roots are the $-\zeta$ with $\zeta$ a primitive $n$-th root (check orders). (c) The left side at $X=1$ is $n$.
\end{hint}
\begin{hint}{ex:napoleon}
With $\alpha=e^{-i\pi/3}=-j$: $A=\frac{(2-\alpha)b+(1+\alpha)c}3$ and similarly for $B,C$ by cyclic permutation. Compute $A+jB+j^2C$ and use $1+j+j^2=0$.
\end{hint}
\begin{hint}{ex:cotsum}
(a) Differentiate $\ln P=\sum\ln(X-\omega^k)$ formally, or use the product rule on $P=\prod(X-\omega^k)$. At $X=1$: $P(1)=n$, $P'(1)=\frac{n(n-1)}2$, $P''(1)=\frac{n(n-1)(n-2)}3$. (b) Half-angle: $1-\omega^k=-2i\sin\frac{k\pi}n\,e^{ik\pi/n}$. Then square and sum, using $\sum\cot\frac{k\pi}n=0$ by the symmetry $k\mapsto n-k$.
\end{hint}

\section{Solutions}

\begin{solution}{ex:lin}
\idea ``Linearise'' means: write a power of $\sin$ as a combination of $\cos(k\theta)$, which are easy to integrate. Either expand $\left(\frac{e^{i\theta}-e^{-i\theta}}{2i}\right)^4$ with the binomial formula, or apply the double-angle formula $\sin^2=\frac{1-\cos2\theta}2$ twice, which is quicker here.
\writeup $\sin^2\theta=\frac{1-\cos2\theta}2$, so
\[
\begin{aligned}\sin^4\theta&=\frac{1-2\cos2\theta+\cos^22\theta}4\eqbr=\frac{1-2\cos2\theta+\frac{1+\cos4\theta}2}4\eqbr=\frac{3-4\cos2\theta+\cos4\theta}8.\end{aligned}
\]
The cosines integrate to $0$ over $[0,\pi]$ (for $k\ge1$, $\int_0^\pi\cos 2k\theta\dd\theta=\left[\frac{\sin2k\theta}{2k}\right]_0^\pi=0$), so $\int_0^\pi\sin^4\theta\dd\theta=\frac{3\pi}8$.
\end{solution}

\begin{solution}{ex:binomcos}
\idea A sum $\sum\binom nk\cos(k\theta)$ looks like the binomial formula with $x=e^{i\theta}$, except that $\cos$ is the real part. So compute the complex sum $\sum\binom nke^{ik\theta}=(1+e^{i\theta})^n$, then use the half-angle factorisation to read off its real part.
\writeup The sum is the real part of $\sum_k\binom nke^{ik\theta}=\left(1+e^{i\theta}\right)^n=\left(2\cos\frac\theta2\right)^ne^{in\theta/2}$, so
\[
\sum_{k=0}^n\binom nk\cos(k\theta)=2^n\cos^n\frac\theta2\,\cos\frac{n\theta}2.
\]
\end{solution}

\begin{solution}{ex:dirichlet}
\idea $D_n$ is a geometric sum with ratio $e^{i\theta}\ne1$. The answer we are given contains half-angles, $\sin\frac\theta2$ and $(n+\frac12)\theta$: this is the signal to use the ``pull out the half-angle'' reflex, i.e.\ to make the numerator and denominator symmetric.
\writeup Since $e^{i\theta}\ne1$, the geometric sum gives
\[
D_n(\theta)=e^{-in\theta}\,\frac{e^{i(2n+1)\theta}-1}{e^{i\theta}-1}.
\]
Multiplying numerator and denominator by $e^{-i\theta/2}$,
\[
D_n(\theta)=\frac{e^{i(n+\frac12)\theta}-e^{-i(n+\frac12)\theta}}{e^{i\theta/2}-e^{-i\theta/2}}
=\frac{\sin\left((n+\frac12)\theta\right)}{\sin\frac\theta2}.
\]
Pairing $k$ with $-k$ gives $D_n(\theta)=1+2\sum_{k=1}^n\cos k\theta$, hence
\[
\sum_{k=1}^n\cos k\theta
=\frac12\left(\frac{\sin\left((n+\frac12)\theta\right)}{\sin\frac\theta2}-1\right).
\]
Figure~\ref{fig:dirichlet} shows $D_n$: a tall peak of height $2n+1$ at $0$, and small oscillations elsewhere. Fourier series (Volume~4) are built on this picture.
\end{solution}

\begin{figure}[ht]
\centering
\fig{ch02-dirichlet}
\caption{The Dirichlet kernel $D_n(\theta)=\dfrac{\sin\left((n+\frac12)\theta\right)}{\sin\frac\theta2}$ for $n=3$ and $n=8$. The peak at $0$ has height $2n+1$ and its first zeros are at $\pm\frac{2\pi}{2n+1}$, so its width is about $\frac{4\pi}{2n+1}$; the area under the curve over $[-\pi,\pi]$ is $2\pi$ for every $n$, since every $e^{ik\theta}$ with $k\ne0$ integrates to $0$.}
\label{fig:dirichlet}
\end{figure}

\begin{solution}{ex:mod3}
\idea Translate: the answer is $\sum_{3\mid k}\binom nk$, the sum of the coefficients of $(1+X)^n$ whose index is in one residue class. That is exactly what the filter of Theorem~\ref{thm:filter} extracts, with $m=3$. The rest is computing $1+j$ in exponential form.
\writeup The number of such subsets is $\sum_{k\equiv0\,(3)}\binom nk$. Apply Theorem~\ref{thm:filter} to $(1+X)^n$ with $m=3$, $r=0$ and
$j=e^{2i\pi/3}$:
\[
\frac13\left(2^n+(1+j)^n+(1+j^2)^n\right).
\]
Now $1+j+j^2=0$ gives $1+j=-j^2=e^{i\pi}e^{-2i\pi/3}=e^{i\pi/3}$ and $1+j^2=\overline{1+j}=e^{-i\pi/3}$, so the answer is $\dfrac{2^n+2\cos\frac{n\pi}3}3$.
For $n=3$ this gives $\frac{8-2}3=2$: the empty set and the whole set.
\end{solution}

\begin{solution}{ex:sines}
\idea A product of sines of angles $\frac{k\pi}n$ suggests the roots of unity, because $|1-e^{i\theta}|=2|\sin\frac\theta2|$ turns $\sin\frac{k\pi}n$ into $\frac12|1-\omega^k|$. So we need $\prod(1-\omega^k)$, which is a value of the polynomial $\prod_{k\ne0}(X-\omega^k)=\frac{X^n-1}{X-1}$.
\writeup Dividing $X^n-1$ by $X-1$ in Proposition~\ref{prop:roots} gives $\prod_{k=1}^{n-1}(X-\omega^k)=1+X+\dots+X^{n-1}$; at $X=1$, $\prod_{k=1}^{n-1}(1-\omega^k)=n$.
Taking moduli, with $|1-e^{i\theta}|=2|\sin\frac\theta2|$ and $\theta=\frac{2k\pi}n$:
\[
n=\prod_{k=1}^{n-1}2\sin\frac{k\pi}n=2^{n-1}\prod_{k=1}^{n-1}\sin\frac{k\pi}n,
\]
all sines being positive since $0<\frac{k\pi}n<\pi$. For $n=7$ this is the $\frac7{64}$ of the opening.
\end{solution}

\begin{solution}{ex:pentagon}
\idea To find two unknown numbers, compute their sum and product: they are then the roots of a quadratic. The relation $1+\omega+\dots+\omega^4=0$ gives the sum, and $\omega^5=1$ lets us reduce exponents in the product. Pairing $\omega$ with $\omega^4=\overline\omega$ makes $c$ real, equal to $2\cos\frac{2\pi}5$.
\writeup By Proposition~\ref{prop:orth}, $1+\omega+\omega^2+\omega^3+\omega^4=0$, so $c+d=-1$. Using $\omega^5=1$,
$cd=\omega^3+\omega^4+\omega^6+\omega^7=\omega^3+\omega^4+\omega+\omega^2=-1$. So $c$ and $d$ are the roots of
$x^2+x-1$, namely $\frac{-1\pm\sqrt5}2$. Since $\omega^4=\overline\omega$, $c=2\cos\frac{2\pi}5>0$ (as $0<\frac{2\pi}5<\frac\pi2$), so
$c=\frac{\sqrt5-1}2$ and $\cos\frac{2\pi}5=\frac{\sqrt5-1}4$.
\end{solution}

\begin{solution}{ex:circulant}
\idea The question hands us candidate eigenvectors; the work in (a) is a change of summation index, which is allowed because indices live modulo $n$ and $\omega^n=1$. For (b), $n$ eigenvectors that form a basis diagonalise $C$, and the determinant is the product of the eigenvalues; the basis property is a Vandermonde determinant. (c) is (b) with $P=1+X$, and $\prod(1+\omega^j)$ is again a value of $X^n-1$, this time at $-1$.
\writeup (a) The $r$-th coordinate of $Cv_j$ is, with $t=s-r$ taken modulo $n$ (allowed since $\omega^{jn}=1$),
\[
\sum_{s=0}^{n-1}a_{(s-r)\bmod n}\,\omega^{js}=\omega^{jr}\sum_{t=0}^{n-1}a_t\,\omega^{jt}=P(\omega^j)\,\omega^{jr}.
\]
(b) The vectors $v_0,\dots,v_{n-1}$ are the columns of a Vandermonde matrix on the distinct numbers $1,\omega,\dots,\omega^{n-1}$; its determinant is non-zero (Chapter~\ref{ch:linear}, Proposition~\ref{prop:vdm}),
so they form a basis. In this basis $C$ is diagonal with entries $P(\omega^j)$, and the determinant, which does not depend on the basis, is their product.

(c) Here $P=1+X$. Evaluating $X^n-1=\prod_j(X-\omega^j)$ at $X=-1$ gives $\prod_j(-1-\omega^j)=(-1)^n-1$, so
\[
\begin{aligned}\det C&=\prod_{j=0}^{n-1}(1+\omega^j)=(-1)^n\left((-1)^n-1\right)\eqbr=1-(-1)^n,\end{aligned}
\]
which is $2$ for $n$ odd and $0$ for $n$ even. (For $n$ even, $(1,-1,1,\dots,-1)$ is in the kernel.)
\end{solution}

\begin{solution}{ex:gauss}
\idea A modulus squared is $G\overline G$, a double sum. Double sums of roots of unity become simple when one variable appears linearly in the exponent, because then orthogonality kills almost everything. The exponent $a^2-b^2$ is quadratic in both, but the substitution $a=b+h$ makes it $h^2+2hb$: linear in $b$.
\writeup Since $\overline\omega=\omega^{-1}$, $|G|^2=G\overline G=\sum_{a,b}\omega^{a^2-b^2}$, where $a,b$ range over $\Z/p\Z$
(the exponent only matters modulo $p$). Substitute $a=b+h$, a bijection for fixed $b$; then $a^2-b^2=h^2+2hb$ and
\[
|G|^2=\sum_{h=0}^{p-1}\omega^{h^2}\sum_{b=0}^{p-1}\omega^{2hb}.
\]
By Proposition~\ref{prop:orth} the inner sum is $p$ if $p\mid2h$ and $0$ otherwise. As $p$ is odd, $p\mid2h$ only for
$h=0$, and $|G|^2=\omega^0\cdot p=p$.
\end{solution}

\begin{solution}{ex:sqrtz}
\idea Square roots of $a+ib$: the system $x^2-y^2=a$, $2xy=b$, completed by $x^2+y^2=|a+ib|$ which comes free from taking moduli. The quadratic equation is then the usual formula, with a complex discriminant whose square root is found the same way.
\writeup For $-7+24i$: $|{-7}+24i|=25$, so $x^2=\frac{25-7}2=9$, $y^2=\frac{25+7}2=16$, and $2xy=24>0$ means $x,y$ have the same sign: the roots are $\pm(3+4i)$. For $2i$: $x^2=y^2=1$, same sign: $\pm(1+i)$.
For $z^2-(1+i)z+i=0$, the discriminant is $(1+i)^2-4i=2i-4i=-2i$, whose square roots are $\pm(1-i)$ (as $(1-i)^2=-2i$). So $z=\frac{(1+i)\pm(1-i)}2$: $z=1$ or $z=i$. Check: $(z-1)(z-i)=z^2-(1+i)z+i$.
\end{solution}

\begin{solution}{ex:sumdist}
\idea Squared distances expand into moduli and real parts of products; the only thing we know about the $z_k$ collectively is that they add up to $0$ (orthogonality with $j=1$, or the centre of mass of a regular polygon), so arrange the computation to make $\sum z_k$ appear.
\writeup Since $|z|=|z_k|=1$, $|z-z_k|^2=(z-z_k)(\bar z-\bar z_k)=|z|^2+|z_k|^2-z\bar z_k-\bar zz_k=2-2\operatorname{Re}(z\bar z_k)$. Summing over $k$,
$\sum_k|z-z_k|^2=2n-2\operatorname{Re}\bigl(z\sum_k\bar z_k\bigr)=2n$, because the $z_k$ are $c\,\omega^k$ for some $c$ with $|c|=1$ and $\sum_k\omega^k=0$ (Proposition~\ref{prop:orth}). For the regular hexagon and $z$ a vertex: $0+1+3+4+3+1=12=2\cdot6$.
\end{solution}

\begin{solution}{ex:cardano}
\idea (a) is pure algebra: $(u+v)^3=u^3+v^3+3uv(u+v)$, which is exactly $x^3=q+px$ if $u^3+v^3=q$ and $3uv=p$. Two numbers with known sum and product are the roots of a quadratic. (b) The formula produces $u^3=2+11i$, and the real question is which cube root to take: there are three, and each pairs with exactly one $v$ through the constraint $uv=5$; the three pairs give the three real roots. This is the point Bombelli could only half see, and Euler cleared up two centuries later.
\writeup (a) $(u+v)^3=u^3+3u^2v+3uv^2+v^3=(u^3+v^3)+3uv(u+v)=q+p(u+v)$, so $x=u+v$ satisfies $x^3=px+q$. The numbers $u^3,v^3$ have sum $q$ and product $(uv)^3=\frac{p^3}{27}$, so they are the roots of $T^2-qT+\frac{p^3}{27}$, namely $\frac q2\pm\sqrt{\frac{q^2}4-\frac{p^3}{27}}$; taking cube roots gives Cardano's formula.

(b) With $p=15$, $q=4$: $\frac{q^2}4-\frac{p^3}{27}=4-125=-121$, so $u^3=2+11i$ and $v^3=2-11i$. Since $(2+i)^3=2+11i$ (checked above), the three cube roots of $2+11i$ are $u_k=(2+i)\omega^k$, $k=0,1,2$, with $\omega=e^{2i\pi/3}=\frac{-1+i\sqrt3}2$ (if $u^3=u_0^3$ then $(u/u_0)^3=1$). The constraint $uv=\frac p3=5$ gives $v_k=\frac5{u_k}=\frac{5}{(2+i)\omega^k}=(2-i)\omega^{-k}$, using $\frac5{2+i}=2-i$. Hence
\[
x_k=u_k+v_k=(2+i)\omega^k+\overline{(2+i)\omega^k}=2\operatorname{Re}\bigl((2+i)\omega^k\bigr).
\]
For $k=0$: $x_0=4$. For $k=1$: $(2+i)\omega=(2+i)\frac{-1+i\sqrt3}2=\frac{-2-\sqrt3+i(2\sqrt3-1)}2$, so $x_1=-2-\sqrt3$. For $k=2$, $\omega^2=\bar\omega$ gives $x_2=-2+\sqrt3$. Check: $4+(-2-\sqrt3)+(-2+\sqrt3)=0$, the coefficient of $x^2$ in $x^3-15x-4$, as Vi\`ete requires (Chapter~\ref{ch:polynomials}).
\end{solution}

\begin{solution}{ex:equi}
\idea (a) The criterion (iii) asks for $a^2+b^2+c^2=ab+bc+ca$; we know $a+b+c=0$, whose square is $a^2+b^2+c^2+2(ab+bc+ca)$, so it suffices to show $ab+bc+ca=0$. On the unit circle, $\bar a=1/a$, which turns $ab+bc+ca$ into $abc(\bar a+\bar b+\bar c)=abc\,\overline{(a+b+c)}=0$. (b) The new triangle is the image of the old one under an explicit linear map; just plug into (ii).
\writeup (a) $0=(a+b+c)^2=a^2+b^2+c^2+2(ab+bc+ca)$ and $ab+bc+ca=abc\left(\frac1c+\frac1a+\frac1b\right)=abc(\bar c+\bar a+\bar b)=abc\,\overline{a+b+c}=0$. So $a^2+b^2+c^2=0=ab+bc+ca$, and $abc$ is equilateral by Proposition~\ref{prop:equilateral}; it is not degenerate since the points are distinct.

(b) If $a+jb+j^2c=0$, then $(a+b)+j(b+c)+j^2(c+a)=a(1+j^2)+b(1+j)+c(j+j^2)=-ja-j^2b-c=-j(a+jb+j^2c)=0$, using $1+j^2=-j$, $1+j=-j^2$, $j+j^2=-1$ and $j^3=1$. The other orientation is the same with $j$ and $j^2$ exchanged. (Geometrically, the new triangle is congruent to the old one: the isometry $z\mapsto3g-z$, with $g=\frac{a+b+c}3$, sends $c,a,b$ to $a+b,b+c,c+a$.)
\end{solution}

\begin{solution}{ex:cyclo}
\idea Everything follows from the factorisation $X^n-1=\prod_{d\mid n}\Phi_d$, used as a recipe: $\Phi_n$ is $X^n-1$ divided by the cyclotomic polynomials of the proper divisors. For (b) and (c), count the roots rather than compute: a cyclotomic polynomial is determined by its roots.
\writeup (a) $\Phi_{10}=\frac{X^{10}-1}{\Phi_1\Phi_2\Phi_5}=\frac{X^5+1}{X+1}=X^4-X^3+X^2-X+1$. Next, $\Phi_{15}=\frac{X^{15}-1}{\Phi_1\Phi_3\Phi_5}=\frac{X^{10}+X^5+1}{X^2+X+1}=X^8-X^7+X^5-X^4+X^3-X+1$ (long division). Finally $\Phi_{18}=\Phi_9(-X)$ by (b), and $\Phi_9=\Phi_3(X^3)=X^6+X^3+1$, so $\Phi_{18}=X^6-X^3+1$.

(b) Let $n$ be odd. If $\zeta$ is a primitive $n$-th root of unity, $-\zeta$ has order $2n$: $(-\zeta)^{2n}=1$, and if $(-\zeta)^m=1$ then $\zeta^{2m}=1$ so $n\mid2m$, hence $n\mid m$ ($n$ odd), and then $(-1)^m=1$ forces $m$ even, so $2n\mid m$. The map $\zeta\mapsto-\zeta$ is injective and $\varphi(2n)=\varphi(n)$ for odd $n$ (Exercise~\ref{ex:phi}; or count: the roots of order $2n$ are the $\omega_{2n}^k$ with $k$ odd and prime to $n$, as many as the $k\le n$ prime to $n$), so the primitive $2n$-th roots are exactly the $-\zeta$. Hence $\Phi_{2n}(X)=\prod_\zeta(X+\zeta)=\prod_\zeta\bigl(-(-X-\zeta)\bigr)=(-1)^{\varphi(n)}\Phi_n(-X)=\Phi_n(-X)$, since $\varphi(n)$ is even for $n\ge3$ (the primitive roots come in conjugate pairs).

(c) At $X=1$, $\frac{X^n-1}{X-1}=1+X+\dots+X^{n-1}$ equals $n$, so $n=\prod_{d\mid n,\,d>1}\Phi_d(1)$. For $n=p$ prime this gives $\Phi_p(1)=p$. For $n=p^k$, the divisors $d>1$ are $p,\dots,p^k$ and induction gives $p^k=p^{k-1}\Phi_{p^k}(1)$, so $\Phi_{p^k}(1)=p$. If $n$ has at least two prime factors, argue by strong induction on $n$: every proper divisor $d>1$ with two prime factors has $\Phi_d(1)=1$ by the induction hypothesis, and the prime-power divisors contribute $\prod_{p^j\mid n}p=\prod_{p\mid n}p^{v_p(n)}=n$. So $n=n\cdot\Phi_n(1)$ and $\Phi_n(1)=1$. Check: $\Phi_6(1)=1-1+1=1$, $\Phi_{12}(1)=1$, $\Phi_{10}(1)=1$.
\end{solution}

\begin{solution}{ex:napoleon}
\idea Everything is linear in $a,b,c$, so the centres are explicit linear combinations, and the criterion of Proposition~\ref{prop:equilateral} is a linear identity in $a,b,c$: it holds for all triangles if and only if the three coefficients vanish, which is a computation in the ring generated by $j$. Since the construction commutes with similarities, we could even assume $a=0$, $b=1$; but the symmetric computation is just as short.
\writeup Let $\alpha=e^{-i\pi/3}=-j$ (indeed $-j=e^{i\pi}e^{2i\pi/3}=e^{5i\pi/3}=e^{-i\pi/3}$). For $abc$ counterclockwise, the outward equilateral triangle on $bc$ has third vertex $p=b+(c-b)\alpha$, and its centre is
\[
A=\frac{b+c+p}3=\frac{(2-\alpha)b+(1+\alpha)c}3=\frac{(2+j)b+(1-j)c}3.
\]
Cyclically, $B=\frac{(2+j)c+(1-j)a}3$ and $C=\frac{(2+j)a+(1-j)b}3$. Then $3(A+jB+j^2C)$ has coefficient of $a$: $j(1-j)+j^2(2+j)=j-j^2+2j^2+1=1+j+j^2=0$; of $b$: $(2+j)+j^2(1-j)=2+j+j^2-1=0$; of $c$: $(1-j)+j(2+j)=1+j+j^2=0$. So $A+jB+j^2C=0$ and $ABC$ is equilateral (counterclockwise) by Proposition~\ref{prop:equilateral}, for every triangle $abc$, even a degenerate one. The theorem is traditionally attributed to Napoleon Bonaparte, who did like mathematics; the earliest known statement is in a problem published by W.~Rutherford in 1825, four years after Napoleon's death, and no evidence links the two.
\end{solution}

\begin{solution}{ex:cotsum}
\idea The logarithmic derivative $\frac{P'}P$ of a product is the sum of the logarithmic derivatives of the factors; this is the product rule, nothing more, and it turns the sum $\sum\frac1{X-\omega^k}$ into something computable from $P$. Differentiating once more produces the sum of the squares. The coefficients of $P$ are all $1$, so $P(1)$, $P'(1)$ and $P''(1)$ are sums of $1$, of $k$ and of $k(k-1)$. Then the half-angle trick converts $\frac1{1-\omega^k}$ into a cotangent, and separating real and imaginary parts gives the sums.
\writeup (a) By the product rule, $P'=\sum_k\prod_{l\ne k}(X-\omega^l)$, so $\frac{P'}P=\sum_k\frac1{X-\omega^k}$ (away from the roots). Differentiating, $\frac{P''P-P'^2}{P^2}=-\sum_k\frac1{(X-\omega^k)^2}$, which is the second formula. At $X=1$: $P(1)=n$, $P'(1)=\sum_{k=1}^{n-1}k=\frac{n(n-1)}2$, $P''(1)=\sum_{k=2}^{n-1}k(k-1)=\frac{(n-1)n(n-2)}3$ (telescoping, Chapter~\ref{ch:sums}). So
\[
\begin{gathered}\sum_{k=1}^{n-1}\frac1{1-\omega^k}=\frac{n-1}2,\\ \sum_{k=1}^{n-1}\frac1{(1-\omega^k)^2}=\frac{(n-1)^2}4-\frac{(n-1)(n-2)}3=\frac{(n-1)(5-n)}{12}.\end{gathered}
\]
(b) With $\theta_k=\frac{k\pi}n$, $1-\omega^k=-2i\sin\theta_k\,e^{i\theta_k}$, so $\frac1{1-\omega^k}=\frac{e^{-i\theta_k}}{-2i\sin\theta_k}=\frac{i(\cos\theta_k-i\sin\theta_k)}{2\sin\theta_k}=\frac12+\frac i2\cot\theta_k$. (The real part $\frac12$ matches the first sum in (a): $\frac{n-1}2$.) Squaring, $\frac1{(1-\omega^k)^2}=\frac14-\frac14\cot^2\theta_k+\frac i2\cot\theta_k$. Sum over $k$: the imaginary parts cancel in pairs since $\cot\theta_{n-k}=-\cot\theta_k$, and comparing with (a),
\[
\begin{gathered}\frac{n-1}4-\frac14\sum_k\cot^2\theta_k=\frac{(n-1)(5-n)}{12},\\ \text{so}\qquad\sum_{k=1}^{n-1}\cot^2\frac{k\pi}n=(n-1)-\frac{(n-1)(5-n)}3=\frac{(n-1)(n-2)}3.\end{gathered}
\]
Adding $n-1$ (as $\frac1{\sin^2}=1+\cot^2$): $\sum_k\frac1{\sin^2\theta_k}=\frac{(n-1)(n-2)}3+(n-1)=\frac{(n-1)(n+1)}3=\frac{n^2-1}3$.

(c) $|1-\omega^k|^2=4\sin^2\theta_k$, so $\sum_k\frac1{|1-\omega^k|^2}=\frac14\cdot\frac{n^2-1}3=\frac{n^2-1}{12}$. For $n=7$: $\frac{48}{12}=4$, which numerical evaluation confirms. As $n\to\infty$ with $k$ fixed, $\sin\frac{k\pi}n\sim\frac{k\pi}n$, so $\frac1{n^2\sin^2(k\pi/n)}\to\frac1{k^2\pi^2}$; dividing the identity by $n^2$ suggests $2\sum_{k\ge1}\frac1{\pi^2k^2}=\frac13$ (the terms with $k$ near $n$ mirror those with $k$ small), i.e.\ $\sum\frac1{k^2}=\frac{\pi^2}6$. Making the exchange of limit and sum rigorous is exactly what Exercise~\ref{ex:basel} does.
\end{solution}

\subsection*{Notes on the challenges}

(1) The values are $G=i\sqrt3$, $\sqrt5$, $i\sqrt7$, $i\sqrt{11}$, $\sqrt{13}$. Gauss recorded the conjecture in his diary in May 1801 and the proof on 30 August 1805, writing that it came ``like a flash of lightning''; the proof is in his \emph{Summatio quarumdam serierum singularium} (1811), and Volume~4 gives one through the Fourier transform.

(2) The admissible $n\le30$ are $3$, $4$, $5$, $6$, $8$, $10$, $12$, $15$, $16$, $17$, $20$, $24$, $30$ (and $n=1,2$ trivially): $\varphi(n)$ must be a power of $2$, which happens when $n=2^a\cdot p_1\cdots p_r$ with distinct Fermat primes $p_i$. The five known Fermat primes are $3,5,17,257,65537$; whether there are more is open. For $n=7$, $\varphi(7)=6$, and $\cos\frac{2\pi}7$ is a root of $8x^3+4x^2-4x-1$, an irreducible cubic: it cannot be reached by square roots.

(3) If $\alpha+\beta=\gamma$ with $|\alpha|=|\beta|=|\gamma|=1$, then dividing by $\gamma$ gives $u+v=1$ with $u,v$ on the unit circle; the rhombus picture (half-angle trick) forces $\{u,v\}=\{e^{i\pi/3},e^{-i\pi/3}\}$. So the solutions are $\gamma\cdot(e^{i\pi/3}+e^{-i\pi/3})$: the triple $(\gamma e^{i\pi/3},\gamma e^{-i\pi/3},\gamma)$, which lives among the $n$-th roots of unity exactly when $6\mid n$. For four roots summing to $0$, the solutions are the pairs of antipodal pairs, $\alpha+(-\alpha)+\gamma+(-\gamma)=0$: this is a theorem of Mann (1965) in its simplest case, and the general classification of vanishing sums of roots of unity is still an active subject.

\section*{Sources and further reading}

\begin{itemize}\raggedright
\item T.~Needham, \emph{Visual Complex Analysis}, 25th anniversary ed., Oxford University Press, 2023, Chapters~1--2. Complex numbers as geometry, with hundreds of pictures.
\item T.~Andreescu and D.~Andrica, \emph{Complex Numbers from A to\dots Z}, 2nd ed., Birkh\"auser, 2014. Roots of unity and many problems like Exercises~\ref{ex:sines}, \ref{ex:pentagon}, \ref{ex:equi} and~\ref{ex:napoleon}.
\item R.~L.~Graham, D.~E.~Knuth and O.~Patashnik, \emph{Concrete Mathematics}, 2nd ed., Addison-Wesley, 1994, Chapter~7. Generating functions and the roots-of-unity filter.
\item T.~H.~Cormen, C.~E.~Leiserson, R.~L.~Rivest and C.~Stein, \emph{Introduction to Algorithms}, 4th ed., MIT Press, 2022, Chapter~30. Roots of unity at work: the fast Fourier transform and fast polynomial multiplication.
\item B.~C.~Berndt, R.~J.~Evans and K.~S.~Williams, \emph{Gauss and Jacobi Sums}, Wiley, 1998. The full story of the sum $G$ of Exercise~\ref{ex:gauss}, including the determination of its sign and Gauss's diary entries about it.
\item P.~J.~Nahin, \emph{An Imaginary Tale: The Story of $\sqrt{-1}$}, Princeton University Press, 1998. Cardano, Bombelli, Wessel (1797, published 1799, rediscovered 1895, French translation 1897), Argand (1806) and Gauss (1831), with the sources.
\item G.~Cardano, \emph{Ars Magna, or The Rules of Algebra}, translated by T.~R.~Witmer, MIT Press, 1968 (Dover reprint 1993). Chapter~37 contains the problem of the two numbers with sum $10$ and product $40$ and the ``mental tortures''.
\item R.~Bombelli, \emph{L'Algebra}, Bologna, 1572; the computation with $x^3=15x+4$ and the ``more sophistic than true'' remark are discussed in J.~Stillwell, \emph{Mathematics and Its History}, 3rd ed., Springer, 2010, Chapter~14, and in Nahin's book.
\item F.~Cajori, \emph{A History of Mathematical Notations}, vol.~2, Open Court, 1929, \S498. Euler's introduction of the symbol $i$ in a memoir presented in 1777 and printed in 1794.
\item J.~J.~Gray, \emph{A commentary on Gauss's mathematical diary, 1796--1814}, Expositiones Mathematicae 2 (1984), 97--130. Entry~1, 30 March 1796: the $17$-gon; entries of 1801 and 1805 on the Gauss sum.
\item A.~Migotti, \emph{Zur Theorie der Kreisteilungsgleichung}, Sitzungsberichte der Wiener Akademie 87 (1883), 7--14. The first cyclotomic polynomial with a coefficient other than $0,\pm1$, $\Phi_{105}$.
\item H.~S.~M.~Coxeter and S.~L.~Greitzer, \emph{Geometry Revisited}, Mathematical Association of America, 1967, Section~3.3. Napoleon's theorem, its history (Rutherford's problem of 1825) and its generalisations.
\item H.~B.~Mann, \emph{On linear relations between roots of unity}, Mathematika 12 (1965), 107--117. Vanishing sums of roots of unity, the third challenge.
\end{itemize}
